\chapter[Band 3: Mengenlehre]{Band 3 auf einen Blick}
\noindent\textbf{Mengenlehre}\par\medskip
\noindent\emph{Lesart.} Alle Gegenstände sind Mengen. \(P\) und
\(\varphi\) bezeichnen Formeln der Mengensprache; bei den Axiomenschemata
gelten die Frischebedingungen des Originals. Die Relations- und
Operatorbegriffe gehören ausdrücklich zu diesem Band.
\begin{BandUebersicht}
\textbf{ZFC und Extensionalität}
Die Mengenstruktur wird durch die ZFC-Axiome beschrieben.
\UebFormel{\begin{aligned}
&\forall x\,(x\in A\leftrightarrow x\in B)\\
&\hspace{15mm}\vdash A=B.
\end{aligned}}
\UebQuellen{\FormulaRefAuto{Menge}[def];
\FormulaRefAuto{\forall x\, (x \in A \leftrightarrow x \in B) \vdash A = B}[ax]}
&
\textit{Gleichheit wird durch Elemente bestimmt}
\UebFormel{\begin{aligned}
&A=B\eqvdash\forall x\,(x\in A\leftrightarrow x\in B),\\
&x\in A\dsep y\notin A\vdash x\neq y.
\end{aligned}}
\UebQuellen{\FormulaRefAuto{\forall x\, (x \in A \leftrightarrow x \in B) \eqvdash A = B}[thm];
\FormulaRefAuto{x\in A\dsep y\notin A \vdash x\neq y}[thm]}\\
\addlinespace[.8ex]
\textbf{Teilmengen und leere Menge}
\UebFormel{\begin{aligned}
A\subseteq B&\coloneq\forall x\,(x\in A\rightarrow x\in B),\\
A\subset B&\coloneq A\subseteq B\land A\neq B,\\
\varnothing&\coloneq\iota O\,\forall x\,(x\notin O).
\end{aligned}}
\UebQuellen{\FormulaRefAuto{A \subseteq B \coloneq \forall x\,(x\in A \rightarrow x\in B)}[def];
\FormulaRefAuto{A \subset B \coloneq A \subseteq B \land A\neq B}[def];
\FormulaRefAuto{\varnothing \coloneq \iota O\bigl(\forall x\,(x \not\in O)\bigr)}[def]}
&
\textit{Leere Menge und echte Erweiterungen}
\UebFormel{\begin{aligned}
&\varnothing\subseteq A,\qquad
A\subseteq\varnothing\vdash A=\varnothing,\\
&A\neq\varnothing\eqvdash\exists x\,(x\in A),\\
&A\subset B\vdash\exists x\in B\setminus A.
\end{aligned}}
\UebQuellen{\FormulaRefAuto{\varnothing\subseteq A}[thm];
\FormulaRefAuto{A\subseteq\varnothing\vdash A=\varnothing}[thm];
\FormulaRefAuto{A\neq\varnothing \eqvdash \exists x\,(x\in A)}[thm];
\FormulaRefAuto{A\subset B\vdash \exists x\in B\setminus A}[thm]}\\
\addlinespace[.8ex]
\textbf{Aussonderung}
Das Aussonderungsaxiom garantiert die Menge der Elemente von \(A\),
die \(P\) erfüllen:
\UebFormel{\{x\in A\mid P(x)\}.}
\UebQuellen{\FormulaRefAuto{\exists B\;\forall x\;\bigl(x \in B \;\leftrightarrow\; x \in A \,\land\, P(x)\bigr)}[ax];
\FormulaRefAuto{\{x\in A \,\mid\, P(x)\} \coloneq \iota B\bigl(\forall x\,(x\in B \leftrightarrow (x\in A \land P(x)))\bigr)}[def]}
&
\textit{Elementkriterium und Grenze der Komprehension}
\UebFormel{\begin{aligned}
&x\in\{u\in A\mid P(u)\}\\
&\hspace{8mm}\eqvdash x\in A\land P(x),\\
&\exists x\,(x\notin A).
\end{aligned}}
Insbesondere gibt es keine Menge aller Mengen.
\UebQuellen{\FormulaRefAuto{x \in \{u \in A \mid P(u)\} \eqvdash x \in A \land P(x)}[thm];
\FormulaRefAuto{FreshElementOutsideSet}[thm]}\\
\addlinespace[.8ex]
\textbf{Schnitt und Differenz}
\UebFormel{\begin{aligned}
A\cap B&\coloneq\{x\in A\mid x\in B\},\\
A\setminus B&\coloneq\{x\in A\mid x\notin B\}.
\end{aligned}}
\UebQuellen{\FormulaRefAuto{A \cap B \coloneq \{ x \in A \mid x \in B \}}[def];
\FormulaRefAuto{A \setminus B \coloneq \{ x \in A \mid x \notin B \}}[def]}
&
\textit{Relative Komplemente}
Für \(X,Y\subseteq M\):
\UebFormel{\begin{aligned}
&M\setminus(M\setminus X)=X,\\
&X\subseteq Y\eqvdash M\setminus Y\subseteq M\setminus X.
\end{aligned}}
\UebQuellen{\FormulaRefAuto{RelCompInvolutive}[thm];
\FormulaRefAuto{RelCompAntitone}[thm]}\\
\addlinespace[.8ex]
\textbf{Paare und geordnete Paare}
\UebFormel{\begin{aligned}
\{a\}&\coloneq\{a,a\},\\
(a,b)&\coloneq\{\{a\},\{a,b\}\}.
\end{aligned}}
Das Paarmengenaxiom liefert \(\{a,b\}\).
\UebQuellen{\FormulaRefAuto{\{a\} \coloneq \{a,a\}}[def];
\FormulaRefAuto{(a, b) \coloneq \{ \{ a \}, \{ a, b \} \}}[def]}
&
\textit{Elemente einer Paarmenge und feste Paar-Tags}
\UebFormel{\begin{aligned}
&x\in\{A,B\}\eqvdash x=A\lor x=B,\\
&(p,x)=(p,y)\vdash x=y,\\
&p\neq q\vdash(p,x)\neq(q,y).
\end{aligned}}
\UebQuellen{\FormulaRefAuto{x \in \{A,B\}\;\eqvdash\;(x = A \lor x = B)}[thm];
\FormulaRefAuto{TaggedPairRightInjective}[thm];
\FormulaRefAuto{TaggedPointsDifferent}[thm]}\\
\addlinespace[.8ex]
\textbf{Vereinigung und Familiendurchschnitt}
\UebFormel{A\cup B\coloneq\bigcup\{A,B\}.}
\(\bigcup M\) sammelt die Elemente der Familienglieder;
\(\bigcap M\) wird für \(M\neq\varnothing\) definiert.
\UebQuellen{\FormulaRefAuto{A \cup B \coloneq \bigcup \{A, B\}}[def];
\FormulaRefAuto{M\neq\varnothing \vdash \bigcap M \coloneq \bigcap\{X\mid X\in M\}}[def]}
&
\textit{Distributivität und gemeinsamer Inhalt}
\UebFormel{\begin{aligned}
&A\cap(B\cup C)=(A\cap B)\cup(A\cap C),\\
&M\neq\varnothing\ \vdash\\
&\quad x\in\bigcap M\leftrightarrow\forall X\in M\,(x\in X).
\end{aligned}}
\UebQuellen{\FormulaRefAuto{A \cap (B \cup C) = (A \cap B) \cup (A \cap C)}[thm];
\FormulaRefAuto{M\neq\varnothing \vdash x\in\bigcap M \leftrightarrow \forall X\in M\,(x\in X)}[thm]}\\
\addlinespace[.8ex]
\textbf{Mengenfamilien}
\UebFormel{\begin{aligned}
&\CoverFam(M,A)\coloneqq A\subseteq\bigcup M,\\
&\PartFam(M,A)\coloneqq\\
&\qquad\DisjFam(M)\land A=\bigcup M.
\end{aligned}}
\(\DisjFam(M)\): nichtleere, paarweise disjunkte Glieder.
\(\UnionClosedFam(M)\): mit \(X,Y\in M\) liegt auch \(X\cup Y\) in \(M\).
Das Disjunktheitsaxiom lautet
\UebFormel{X,Y\in M\dsep X\neq Y\vdash X\cap Y=\varnothing.}
\UebQuellen{\FormulaRefAuto{\CoverFam(M,A) \coloneqq A \subseteq \bigcup M}[def];
\FormulaRefAuto{\PartFam(M,A) \coloneqq \DisjFam(M)\land A=\bigcup M}[def];
\FormulaRefAuto{Disjunkte Familie}[def];\FormulaRefAuto{Vereinigungsabgeschlossene Familie}[def]}
&
\textit{Eindeutigkeit des Familiengliedes}
Für \(\DisjFam(M)\):
\UebFormel{\begin{aligned}
&X,Y\in M\dsep a\in X\dsep a\in Y\vdash X=Y,\\
&X,Y\in M\dsep X\neq Y\dsep a\in X\\
&\hspace{10mm}\vdash a\notin Y.
\end{aligned}}
\UebQuellen{\FormulaRefAuto{X \in M \dsep Y \in M \dsep a \in X \dsep a \in Y \vdash X = Y}[thm];
\FormulaRefAuto{X \in M \dsep Y \in M \dsep X \neq Y \dsep a \in X \vdash a \notin Y}[thm]}
\\\addlinespace[.8ex]
\textbf{Regularität}
\UebFormel{\begin{aligned}
&A\neq\varnothing\\
&\hspace{8mm}\vdash\exists x\in A\,(x\cap A=\varnothing).
\end{aligned}}
\UebQuellen{\FormulaRefAuto{A \neq \varnothing \vdash \exists x \in A \,(x \cap A = \varnothing)}[ax]}
&
\textit{Keine Selbstmitgliedschaft}
\UebFormel{\begin{aligned}
&a\notin a,\qquad a\in b\vdash b\notin a,\\
&a\cup\{a\}=b\cup\{b\}\eqvdash a=b.
\end{aligned}}
\UebQuellen{\FormulaRefAuto{a\not\in a}[thm];
\FormulaRefAuto{a\in b\vdash b\not\in a}[thm];
\FormulaRefAuto{a\cup\{a\}=b\cup\{b\}\eqvdash a=b}[thm]}\\
\addlinespace[.8ex]
\textbf{Potenzmenge und Mengensystem}
\(\powerset(A)\) ist die durch das Potenzmengenaxiom bestimmte Menge
aller Teilmengen von \(A\).
\UebFormel{\begin{aligned}
&\mathsf{SetSys}(M;X)\coloneqq M\subseteq\powerset(X),\\
&\mathfrak I[P,Q]\coloneqq\\
&\qquad\{H\in\powerset(Q)\mid P\subseteq H\}.
\end{aligned}}
\UebQuellen{\FormulaRefAuto{SetSystemOverDef}[def];
\FormulaRefAuto{BooleanSetIntervalDef}[def]}
&
\textit{Elementkriterium und Monotonie}
\UebFormel{\begin{aligned}
&x\in\powerset(A)\eqvdash x\subseteq A,\\
&A\subseteq B\vdash\powerset(A)\subseteq\powerset(B),\\
&\powerset(\varnothing)=\{\varnothing\}.
\end{aligned}}
\UebQuellen{\FormulaRefAuto{x \in \powerset(A)\;\eqvdash\; x \subseteq A}[thm];
\FormulaRefAuto{A \subseteq B \vdash \powerset(A) \subseteq \powerset(B)}[thm];
\FormulaRefAuto{\powerset(\varnothing)=\{\varnothing\}}[thm]}\\
\addlinespace[.8ex]
\textbf{Kartesisches Produkt, Relation und Graph}
\(A\times B\) enthält die geordneten Paare aus \(A\) und \(B\).
\UebFormel{\begin{aligned}
&F\subseteq A\times B\quad\text{Relation},\\
&\GraphStruct{V,E}\coloneqq E\subseteq V\times V.
\end{aligned}}
\UebQuellen{\FormulaRefAuto{F \subseteq A \times B}[def];
\FormulaRefAuto{GraphDef}[def]}
&
\textit{Paar- und Produktkriterien}
\UebFormel{\begin{aligned}
&(a,b)\in A\times B\eqvdash a\in A\land b\in B,\\
&A\times\varnothing=\varnothing,\\
&\varnothing\times A=\varnothing.
\end{aligned}}
Eine Relation muss hier weder total noch funktional sein; diese
Axiome folgen erst in den eigenständigen Bänden 4 und 5.
\UebQuellen{\FormulaRefAuto{(a,b)\in A\times B\eqvdash a\in A\land b\in B}[thm];
\FormulaRefAuto{A\times\varnothing=\varnothing}[thm];
\FormulaRefAuto{\varnothing\times A=\varnothing}[thm]}\\
\addlinespace[.8ex]
\textbf{Binärer Operator}
\UebFormel{\mathsf{BinOp}(A,\star).}
Das Strukturaxiom fordert Abgeschlossenheit:
\UebFormel{x,y\in A\vdash x\star y\in A.}
\UebQuellen{\FormulaRefAuto{BinaryOperationDef}[def];
\FormulaRefAuto{BinaryOperationClosureAxiom}[ax]}
&
\textit{Reichweite dieses Grundbegriffs}
Der Band stellt damit die Typisierung algebraischer Terme bereit.
Assoziativität, Kommutativität, Idempotenz und Kürzbarkeit werden
hier noch nicht vorausgesetzt. Die entsprechenden eigenständigen
Strukturbände ergänzen jeweils ihre eigenen Axiome.\\
\addlinespace[.8ex]
\textbf{Ersetzung, Auswahl und Unendlichkeit}
Ersetzung sammelt eindeutig durch eine Formel bestimmte Werte.
Das Auswahlaxiom trifft jedes Glied einer disjunkten nichtleeren
Familie genau einmal. Unendlichkeit fordert eine induktive Menge:
\UebFormel{\exists A\,(\varnothing\in A\land
\forall x\in A\,(x\cup\{x\}\in A)).}
\UebQuellen{\FormulaRefAuto{ReplacementSchema}[ax];
\FormulaRefAuto{ChoiceFamilyAxiom}[ax];
\FormulaRefAuto{\exists A\;\bigl(\varnothing \in A \land \forall x \in A\,(x \cup\{x\} \in A)\bigr)}[ax]}
&
\textit{Durch einen Mengenterm beschriebene Bildmenge}
Für einen Mengenterm \(t(j)\) mit festen Parametern:
\UebFormel{\begin{aligned}
&\exists!X\,\forall x\,\bigl(\\
&\quad x\in X\leftrightarrow
\exists j\,(j\in J\land x=t(j))\bigr).
\end{aligned}}
Dies folgt aus Ersetzung und Extensionalität und setzt noch
keinen Funktionsgraphen voraus.
\UebQuellen{\FormulaRefAuto{TermImageSetUniqueExists}[thm]}\\
\end{BandUebersicht}
